\begin{pblock}{Choosing the water detector}
Four detectors were run in \term{one} backend job over byte-identical dates
and cloud masks, so every difference is the index and not the sampling.
Scored by rank correlation between water frequency and LiDAR elevation —
a near model-free arbiter, since a lower pixel must flood more often:

\vspace{3mm}
\begin{tabular}{@{}l r l@{}}
\toprule
Index & {Spearman $\rho$} & Bands \\
\midrule
\textbf{NDWI}  & -0.624 & B03/B08, native 10\,m \\
SCL (class 6)  & -0.586 & scene classification, 20\,m \\
AWEI           & -0.584 & four-band \\
MNDWI          & -0.559 & B03/B11, SWIR at 20\,m \\
\bottomrule
\end{tabular}

\vspace{5mm}
NDWI wins for a physical reason: both its bands are natively 10\,m,
so no resampling smears the waterline that carries the signal. MNDWI's SWIR
band is the problem — wet sediment reflects like water, so it over-detects.
Over the whole Bay of Santander (220\,km\textsuperscript{2}, 247
scenes) SCL and NDWI agree on 91.7\,\% of pixels, while MNDWI
claims \emphnum{five times} the intertidal area.

\vspace{4mm}
\pfig[\linewidth]{29_santander_tres}
\figcap{Bay of Santander: the three detectors on identical observations, and
the two difference maps. Red is where the other detector sees more water.}
\end{pblock}
